% Options for packages loaded elsewhere
\PassOptionsToPackage{unicode}{hyperref}
\PassOptionsToPackage{hyphens}{url}
\documentclass[
  ignorenonframetext,
]{beamer}
\newif\ifbibliography
\usepackage{pgfpages}
\setbeamertemplate{caption}[numbered]
\setbeamertemplate{caption label separator}{: }
\setbeamercolor{caption name}{fg=normal text.fg}
\beamertemplatenavigationsymbolsempty
% remove section numbering
\setbeamertemplate{part page}{
  \centering
  \begin{beamercolorbox}[sep=16pt,center]{part title}
    \usebeamerfont{part title}\insertpart\par
  \end{beamercolorbox}
}
\setbeamertemplate{section page}{
  \centering
  \begin{beamercolorbox}[sep=12pt,center]{section title}
    \usebeamerfont{section title}\insertsection\par
  \end{beamercolorbox}
}
\setbeamertemplate{subsection page}{
  \centering
  \begin{beamercolorbox}[sep=8pt,center]{subsection title}
    \usebeamerfont{subsection title}\insertsubsection\par
  \end{beamercolorbox}
}
% Prevent slide breaks in the middle of a paragraph
\widowpenalties 1 10000
\raggedbottom
\AtBeginPart{
  \frame{\partpage}
}
\AtBeginSection{
  \ifbibliography
  \else
    \frame{\sectionpage}
  \fi
}
\AtBeginSubsection{
  \frame{\subsectionpage}
}
\usepackage{iftex}
\ifPDFTeX
  \usepackage[T1]{fontenc}
  \usepackage[utf8]{inputenc}
  \usepackage{textcomp} % provide euro and other symbols
\else % if luatex or xetex
  \usepackage{unicode-math} % this also loads fontspec
  \defaultfontfeatures{Scale=MatchLowercase}
  \defaultfontfeatures[\rmfamily]{Ligatures=TeX,Scale=1}
\fi
\usepackage{lmodern}
\ifPDFTeX\else
  % xetex/luatex font selection
\fi
% Use upquote if available, for straight quotes in verbatim environments
\IfFileExists{upquote.sty}{\usepackage{upquote}}{}
\IfFileExists{microtype.sty}{% use microtype if available
  \usepackage[]{microtype}
  \UseMicrotypeSet[protrusion]{basicmath} % disable protrusion for tt fonts
}{}
\makeatletter
\@ifundefined{KOMAClassName}{% if non-KOMA class
  \IfFileExists{parskip.sty}{%
    \usepackage{parskip}
  }{% else
    \setlength{\parindent}{0pt}
    \setlength{\parskip}{6pt plus 2pt minus 1pt}}
}{% if KOMA class
  \KOMAoptions{parskip=half}}
\makeatother
\setlength{\emergencystretch}{3em} % prevent overfull lines
\providecommand{\tightlist}{%
  \setlength{\itemsep}{0pt}\setlength{\parskip}{0pt}}
\usepackage{bookmark}
\IfFileExists{xurl.sty}{\usepackage{xurl}}{} % add URL line breaks if available
\urlstyle{same}
\hypersetup{
  hidelinks,
  pdfcreator={LaTeX via pandoc}}

\author{\texorpdfstring{}{}}
\date{}

\begin{document}

\begin{frame}[fragile]{Abstract}
\protect\phantomsection\label{abstract}
The reproducibility crisis in computational research is fundamentally
structural: research artifacts are scattered across disconnected
tools---LaTeX editors, Jupyter notebooks, ad-hoc shell scripts---with no
enforced mechanism to keep code, data, and manuscript synchronized.
Studies have shown that most published findings are false positives,
replication rates in psychology hover around 36\%, and only 24\% of 1.4
million Jupyter notebooks can be successfully re-executed. Existing
tools address fragments of this problem: workflow managers (Snakemake,
Nextflow, CWL) orchestrate computation; literate programming systems
(Quarto, Jupyter Book, R Markdown, Overleaf, OpenAI Prism) render
documents; data versioning tools (DVC) track artifacts---but none
enforces cross-cutting quality standards as architectural invariants.
\texttt{template/} applies the principle of Infrastructure as Code to
the research lifecycle, making the manuscript, test suite, and
provenance chain version-controlled, deterministically buildable, and
independently verifiable. It is built on a Two-Layer Architecture that
separates 13 reusable infrastructure subpackages (\textasciitilde271
Python modules, validated by \textasciitilde6,021 tests) from
self-contained project workspaces, connected by an ten-stage DAG-based
build pipeline progressing from environment sanitization through test
execution (with a Zero-Mock testing policy enforcing 90\% project-level
and 60\% infrastructure-level coverage via real filesystem operations
and subprocess invocations), analysis script invocation, Pandoc/XeLaTeX
rendering, SHA-256 cryptographic hashing with steganographic
watermarking, structural PDF validation, and LLM-assisted review. A
Documentation Duality standard equips every directory with both
human-readable \texttt{README.md} and machine-readable
\texttt{AGENTS.md} files, while each infrastructure module additionally
carries a \texttt{SKILL.md}---a structured skill descriptor aligned with
the Model Context Protocol---enabling AI agents to locate and invoke
module capabilities without hallucinating API signatures.

Scalability is demonstrated across three heterogeneous projects---a
gradient descent study (\texttt{code\_project}, 39 tests), a
compositional case modeling study (\texttt{cognitive\_case\_diagrams},
912 tests), and this self-referential architectural analysis
(\texttt{template}, 66 tests)---achieving 100\% pipeline success with
zero mock violations. The fact that these words, these metrics, and the
figures accompanying them were generated by the very pipeline they
describe is itself a demonstration of the system's self-productive
capacity: the manuscript is not merely \emph{about} \texttt{template/}
but \emph{of} it, rendered through the same ten-stage DAG pipeline,
validated by the same test suite, and watermarked by the same
steganographic layer documented herein. A comparative feature analysis
against nine peer tools across fourteen dimensions confirms that
\texttt{template/} uniquely integrates all fourteen distinctive
capabilities---testing enforcement, coverage thresholds, cryptographic
provenance, steganographic watermarking, multi-project management,
AI-agent documentation, agentic skill protocol, interactive TUI,
Zero-Mock policy, manuscript rendering, and pipeline
orchestration---within a single enforced pipeline. \texttt{template/} is
open source under the Apache 2.0 License at
\texttt{github.com/docxology/template}, and is presented as a work in
progress.~
\end{frame}

\end{document}
